\section{Results}
\subsection{Duration changes the amount recovered}
The primary severity~2 result is large. Recovery increases from $R=+0.085$ \ci{+0.066}{+0.105} at $3.84$\,s to $+0.244$ \ci{+0.215}{+0.273} at $10.24$\,s, giving $J=+0.159$ \ci{+0.131}{+0.188}. The four-point sweep in Table~\ref{tab:evidence} shows a gradual increase rather than an endpoint anomaly. The effect also has practical scale. Recovery closes $44\%$ of the gap from P to dense at $3.84$\,s and $82\%$ at $10.24$\,s. Relative to real audio, the corresponding fractions are $33\%$ and $63\%$.

The robustness evidence is now visible in the paper rather than delegated to the repository. Human-CLAP, KL to real references and PANNs top-10 capture all show the same larger recovery at longer durations. The last $7.68$ to $10.24$\,s increment is not resolved by the secondary metrics, suggesting that most of their increase has already accrued by $7.68$\,s. The published DDIM $200$, guidance $3.5$ recipe reproduces the endpoint interaction with $J=+0.184$ \ci{+0.126}{+0.243}, and a Holm correction over the registered severity~2 family leaves the conclusions unchanged.

Two follow-ups clarify the shape. At $15.36$\,s, recovery is $+0.264$ \ci{+0.216}{+0.310}. On the matched $96$-prompt subset, the step beyond $10.24$\,s is only $+0.021$ \ci{-0.023}{+0.067}. The gain therefore plateaus rather than peaking sharply at the native duration. Severity~1 also resolves with more data. Adding $96$ prompts to the original $80$ gives a pooled $J=+0.112$ \ci{+0.076}{+0.149}. The new subset has a larger interaction than the original subset, so the pooled value should be read with that between-set heterogeneity in mind.

\subsection{The duration effect is not training-duration specialization}
The most direct mechanism test changes the duration used for fine-tuning. We fine-tuned severity~2 P for $20{,}000$ steps on $3.84$\,s crops. At its own training duration, the resulting gain is only $+0.009$ \ci{-0.006}{+0.024}. Evaluated at $10.24$\,s, the same checkpoint gains $+0.075$ \ci{+0.053}{+0.096}, yielding $J=+0.065$ \ci{+0.043}{+0.087}. If recovery specialized to the duration on which adaptation was performed, this intervention should have made the short-duration gain larger. It produces the opposite ordering.

A second, weaker reference points in the same direction. A publicly released dense AudioLDM-M text-fine-tuned checkpoint changes by $-0.022$ \ci{-0.061}{+0.017} relative to its dense base at $3.84$\,s and by $+0.091$ \ci{+0.042}{+0.141} at $10.24$\,s, with $J=+0.113$ \ci{+0.051}{+0.173}. This checkpoint was trained with a different recipe and corpus, so it is not the missing matched dense control. It does show that a duration-dependent fine-tuning gain is not unique to post-pruning recovery.

\subsection{Domain transfer is selective rather than absent}
The original hip-hop result looked like a failure of transfer, but the expanded domain study changes that interpretation. Recovery transfers strongly to Clotho. At $10.24$\,s, Clotho gives $R=+0.210$ \ci{+0.176}{+0.243} and closes $74\%$ of the dense gap. On a matched $96$-prompt comparison, the difference between AudioCaps and Clotho recovery is only $+0.032$ \ci{-0.023}{+0.088}. A shift away from AudioCaps therefore does not by itself erase recovery.

Hip-hop remains different. After extending the battery to $127$ prompts, the gain becomes small but resolved, $+0.026$ \ci{+0.008}{+0.044} at $3.84$\,s and $+0.027$ \ci{+0.004}{+0.051} at $10.24$\,s. The dense anchors answer the floor objection directly. On the original $64$-prompt native cell, dense lies $+0.106$ \ci{+0.079}{+0.135} above shuffled-caption chance, while recovery closes only about $2\%$ of the dense gap. The battery can discriminate alignment, but the recovered checkpoint captures little of the available improvement. Because hip-hop changes both content and caption style, the correct conclusion is selective transfer, not a general in-domain versus out-of-domain dichotomy.

\subsection{The short-duration deficit arises during generation}
The duration interaction is not explained by where useful events appear in the clip. FineLAP finds the native-duration grounding gain distributed across time, with no early-versus-late difference, $T=-0.002$ \ci{-0.024}{+0.020}. A crop analysis reaches the same conclusion from another direction. Scoring only the first $3.84$\,s of a clip generated at $10.24$\,s gives $R_{\rm crop}=+0.172$ \ci{+0.150}{+0.194}, which is $+0.087$ \ci{+0.065}{+0.110} above recovery from a separately generated $3.84$\,s clip. The deficit therefore arises mainly when the model is asked to generate at the short operating point, not because CLAP sees a short window.

